\paragraph{The kernel witness on all six signed source blocks.}
The element $k$ of B.1.3 can be chosen to satisfy
\[
 f_k,\beta_k,A_k\in F'\cap\operatorname{core}_B\ker\varphi_P.
                                                        \tag{B.M14}
\]
This is a direct strengthening of its already executed
compactness construction.  For every $J$, (B.C21) gives
$f_J,\beta_J\in R_{\rm nil}(N_J)\subseteq N_J$, and
$N_J$ is the actual sixfold core in (B.C18).  Add all six
closed conditions
$\varphi_P(bfb^{-1}),\varphi_P(b\beta b^{-1}),
\varphi_P(bAb^{-1})\in J$
to (B.C22).  The same construction with the intersection
of the finitely many $J$'s satisfies them all.  Compactness
therefore gives (B.M14), with every previous frame and
character value retained.  For the cubic frame, its
retained common $D_J$-value is one, so $A_J\in N_J$;
and $(1-\operatorname{Ad}_r)A_J^{\mathrm{ab}}=0$ with
determinant three gives $A_J\in F'$.

Every signed block $\theta(bf_kb^{-1})$,
$\theta(b\beta_kb^{-1})$ and $\theta(bA_kb^{-1})$
lies in $\{\pm I\}$, since
its image in the actual projective quotient is trivial.
Each also lies in $\operatorname{SL}_2(\mathbb Z/12)'$,
because $F'$ is preserved by the deck conjugations.
The sign $-I$ is excluded already modulo four.  Explicitly,
put
\[
 K_4=\{I,-U^2,-L^2,U^2L^2\},\qquad
 H_4=\bigcup_{j=0}^2 K_4R^j
       \ \leq\operatorname{SL}_2(\mathbb Z/4),
                                                        \tag{B.M15}
\]
where $L=\left(\begin{smallmatrix}1&0\\-1&1\end{smallmatrix}\right)$
and $R$ is the matrix in (B.M1).  Matrix multiplication
gives $|H_4|=12$, and its four cosets $H_4U^j$ are
the entire group of order forty-eight.  The assignment
$H_4U^j\mapsto j\bmod4$ is a homomorphism, with
$U\mapsto1$, $S\mapsto3$ and $-I\mapsto2$.
The exact enumeration in \texttt{check\_signed\_core.py}
checks every one of the $48^2$ products in this stated
finite matrix group.  Hence $-I$ is not in its derived
subgroup.  All six signed blocks in question are
therefore $I$.

Their original modular blocks are also one:
$q_{12}(\beta_k)=1$ by construction and
$q_{12}(f_k)=q_{12}(s\beta_ks\beta_k^{-1})=1$.
Thus both $f_k$ and $\beta_k$ are trivial in the
specific $D$ of (B.M2), so the cusp-fixed action of
$k$ on $D$ is the identity.  Its action on $D^{\mathrm c}$
is likewise the identity by the singleton-fibre
assertion of B.2.6.  Its actual pure cubic frame
therefore gives an element
\[
 c\in C_{D^{\mathrm c}}(R_0)\cap P,
             \qquad q_{\mathrm c}(c)=[5I].      \tag{B.M16}
\]
Its image in $D$ is the specific central order-two
element with six identity signed blocks and final
modular block $[5I]$.  The word $A_0$ of (B.C17)
already represents this element.
This proves the required signed-source assertion;
it is stronger than the outer equality $D_P(k)=1$.

\paragraph{Literal retention under commutator normalization.}
For the finite group $D$ used here, the exact pure
abelianization and two orders are
\[
 \overline F_0^{\mathrm{ab}}\cong C_{12}^2,
 \qquad |[\overline x]-[\overline y]|=12,
 \qquad |\overline x|=12.                       \tag{B.M17}
\]
They are computed with the actual quotient map
$\overline F_0\to\overline F_0/\overline F_0'$
in \path{check_source_abelian_normalization.g}.
For a full cusp-fixed lift with $f_D\in\overline F_0'$,
the exponent $a$ in (B.M13) therefore belongs to
$12\widehat{\mathbb Z}$.  Thus $\overline x^a=1$,
and commutator normalization preserves its literal
$D$-action and both $D$-frame residues.
